\documentclass[11pt]{article}
\usepackage[T1]{fontenc}
\usepackage{xcolor}
\usepackage{xurl}
\usepackage{lmodern,amsmath,amssymb,amsthm,mathtools,booktabs,geometry,microtype,hyperref}
\geometry{margin=1.05in}
\hypersetup{colorlinks=true,linkcolor=blue!50!black,urlcolor=blue!60!black,citecolor=blue!50!black,pdftitle={The sharp localization barrier for positive Dirichlet zero filters},pdfsubject={Entire division and endpoint rigidity}}
\newtheorem{theorem}{Theorem}[section]
\newtheorem{lemma}[theorem]{Lemma}
\newtheorem{corollary}[theorem]{Corollary}
\newtheorem{proposition}[theorem]{Proposition}
\theoremstyle{definition}\newtheorem{definition}[theorem]{Definition}
\theoremstyle{remark}\newtheorem{remark}[theorem]{Remark}
\newcommand{\R}{\mathbb R}\newcommand{\C}{\mathbb C}\newcommand{\dd}{\,\mathrm d}\newcommand{\ord}{\operatorname{ord}}\newcommand{\Rea}{\operatorname{Re}}\newcommand{\Ima}{\operatorname{Im}}\newcommand{\F}{\mathcal F}\newcommand{\G}{\mathcal G}\newcommand{\calL}{\mathcal L}
\title{\textbf{The sharp localization barrier for positive Dirichlet zero filters}\\[5pt]\large A factor-of-two improvement and endpoint rigidity via entire division}
\author{Research note for the ProveIt GRH project}
\date{27 November 2026}
\begin{document}
\maketitle
\begin{abstract}
The merged ProveIt article asks whether the bound $b\le 2\pi d/Q_*$ is sharp for a positive probability measure cancelling all noncentral zeros, with multiplicities, of a conjugation-closed family of $d$ primitive Dirichlet $L$-functions and possessing a finite moment $\int\exp(b\exp(2|t|/d))\dd\mu(t)$. Assuming the Dirichlet GRH input of that article, we prove the strictly stronger universal bound $b\le \pi d/Q_*$. At the new critical value, the bilateral Laplace transform of every such measure must be a polynomial multiple of the normalized completed product with its central zero removed. The degree is constrained by the gamma-factor power and the central order; positivity of a finite-dimensional family of differential polynomials is a necessary and sufficient test for attainment. For any product of zeta, Dirichlet beta, and the real character modulo five, the new bound is the exact supremum, and an explicit inequality decides attainment. The canonical positive kernel also gives the exact constant for its double-exponential tail and for the logarithmic radius needed to truncate a global averaging rule. No simplicity or independence of ordinates is assumed. The result answers the repository's sharpness question in the negative and settles the revised extremum in an infinite family; it does not assert a solution to an unrelated named conjecture.
\end{abstract}
\tableofcontents

\section{Source, problem and principal results}
\paragraph{Repository analysis.} We inspected the project README, the merged article's README and introduction, theorem headings and open-question lists across its analytic and transfer sections, and the global-filter construction in \texttt{Article/sections/fil.tex} \cite{proveit,fil}. The article merges 209 reports (982 pages according to its README), explicitly does not reverify them, and treats the Dirichlet GRH theorem exported by \texttt{lsernau/grh-lean} as an input. The vendored upstream README pins commit \texttt{72fead102f663d5c7a407964a13dc0cf9d977207}; \texttt{src/Statements.lean} states the relevant strip and half-plane properties in Mathlib terms and \texttt{src/Main.lean} exports them \cite{lean}. Its README reports independent checker acceptance; the vendored package itself is described as unbuilt in ProveIt. We inspected the statements and wrappers but did not rebuild or audit the entire upstream proof. Accordingly the results here are stated as implications from Dirichlet GRH, not as an independently established proof of GRH.

The article's transferable pattern compares two estimates on a common quantity and uses analytic growth to constrain zero locations. In its section on positive filters it already constructs full-spectrum filters with optimal exponent $2/d$ and obtains an endpoint coefficient from circular Jensen averaging. The specific open question at the end of that section asks whether $2\pi d/Q_*$ is sharp (labels \texttt{fil:cor:endpoint}, \texttt{fil:sec:open}). We replace averaging of growth over a circle by division by the \emph{entire prescribed divisor}. The gamma factors then constrain growth in a wedge around the real axis, so additional zeros cannot improve localization beyond a sharper barrier. We give the proof independently, using standard facts about Dirichlet $L$-functions and entire functions.

\subsection{Notation and hypotheses}
Let $\F$ be a nonempty finite multiset of primitive Dirichlet characters, closed under conjugation; repetitions count. The conductor-one character denotes $\zeta$. Put
\[
 d=|\F|,\quad q_\chi=\text{conductor of }\chi,\quad Q_*=(\prod_{\chi\in\F}q_\chi)^{1/d},
 \quad A=d/2,\quad \alpha=1/A=2/d,\quad b_0=\pi d/Q_*.
\]
For a nonprincipal primitive character put $a_\chi=(1-\chi(-1))/2$ and
\[
 \Lambda_\chi(s)=(q_\chi/\pi)^{(s+a_\chi)/2}\Gamma((s+a_\chi)/2)L(s,\chi).
\]
For the conductor-one character put $a_\chi=0$, $\Lambda_\chi(s)=\xi(s)=\tfrac12s(s-1)\pi^{-s/2}\Gamma(s/2)\zeta(s)$. Let
\[
 G(w)=\prod_{\chi\in\F}\Lambda_\chi(\tfrac12+w),\qquad m=\ord_{w=0}G(w),\qquad
 P(w)=\frac{G(w)/w^m}{(G(w)/w^m)|_{w=0}}.
\]
The value at zero is by analytic continuation, so $P(0)=1$. Write $r_0$ for the number of conductor-one factors, $r_1$ for the number of odd factors, and define the gamma power
\begin{equation}\label{eq:kappa}
 \kappa=2r_0+\frac{2r_1-d}{4}-m.
\end{equation}
Under Dirichlet GRH, $P$ is a real even entire function of order at most one; its zero multiset consists of $i\gamma$ at the noncentral ordinates of the family, with aggregate multiplicities $m_\gamma$. The functional equations, closure under conjugation and removal of the central factor give evenness (the sign of $G(-w)/G(w)$ is $(-1)^m$). The usual nonvanishing on the edges of the critical strip and the cancellation of trivial zeros in the completions ensure there are no other zeros. In particular $P(u)>0$ on the real axis.

\begin{definition} A \emph{full-spectrum filter} is a Borel probability measure $\mu$ on $\R$ whose bilateral Laplace transform
\[
 F(w)=\int_\R e^{wt}\dd\mu(t)
\]
is entire and has $\ord_{i\gamma}F\ge m_\gamma$ for every $\gamma\ne0$. We call $b>0$ admissible for $\mu$ if
\begin{equation}\label{eq:moment}
 M_b(\mu)=\int_\R \exp\bigl(b e^{\alpha |t|}\bigr)\dd\mu(t)<\infty.
\end{equation}
In our use, the moment condition guarantees the entirety of the integral. We allow nonsymmetric measures; central zeros are not required to vanish because $F(0)=1$.
\end{definition}

\begin{theorem}[Localization barrier and endpoint rigidity]\label{thm:main}
Assume Dirichlet GRH. For any full-spectrum filter and any admissible $b$, one has
\begin{equation}\label{eq:main} b\le b_0=\pi d/Q_* .\end{equation}
If $b=b_0$ is admissible, there is a real polynomial $Q$, with $Q(0)=1$, $Q(u)>0$ for all $u\in\R$, such that
\begin{equation}\label{eq:rigid}F(w)=P(w)Q(w),\qquad \deg Q+\kappa<-\tfrac12.\end{equation}
If $\mu$ is even then $Q$ is even. In particular the endpoint is impossible when $\kappa\ge-1/2$.
\end{theorem}

\begin{theorem}[Exact endpoint criterion]\label{thm:criterion}
Define the smooth real even signed kernel
\[
 h(t)=\frac1{2\pi}\int_\R P(iv)e^{-ivt}\dd v.
\]
It and all its derivatives decay faster than every ordinary exponential. The endpoint $b_0$ is admissible for some full-spectrum filter if and only if there exists a real polynomial $Q$ with $Q(0)=1$, $\deg Q+\kappa<-1/2$, for which $Q(-\partial_t)h(t)\ge0$ for all $t\in\R$. All endpoint filters are precisely the measures with these densities. For existence one can restrict to even $Q$ by symmetrizing. If $h$ is nonnegative, the optimal supremum is $b_0$ and the canonical filter $h(t)\dd t$ attains it exactly when $\kappa<-1/2$.
\end{theorem}

\paragraph{Concrete families.} For $r$ copies of $\zeta$, $s$ copies of $L(s,\chi_{-4})$, and $j$ copies of $L(s,\chi_5)$, where $\chi_5$ is the real quadratic character modulo five, $d=r+s+j>0$ and
\begin{equation}\label{eq:example} b_0=\frac{\pi(r+s+j)}{(4^s5^j)^{1/d}},\qquad \kappa=\frac{7r+s-j}{4}.\end{equation}
The supremum is $b_0$ and is attained if and only if $j>7r+s+2$. Thus $\zeta$ alone has supremum $\pi$, not $2\pi$, but no filter has a finite moment at $b=\pi$. Three copies of $\chi_5$ have supremum $3\pi/5$ and do attain it. The proof includes the positivity of the kernels needed in these examples.

\section{The analytic estimates}
We use classical analytic continuation, the functional equation, order-one Hadamard factorization, Euler products on the right, the Riemann--von Mangoldt zero count with $O(\log T)$ error for fixed conductor, uniform complex Stirling, and the Phragm\'en--Lindel\"of principle. Standard references are \cite{davenport,titch,zeros,boas,hormander}. All implicit constants in this note may depend on the fixed family.

\begin{lemma}[Gamma growth in the right half-plane]\label{lem:stirling}
For $w=u+iv$, $u\ge2$, with the principal logarithm, uniformly as $|w|\to\infty$,
\begin{equation}\label{eq:stirling}
 \log |P(w)|=A\Rea\!\left(w\log w+(\log(Q_*/(2\pi))-1)w\right)+\kappa\log|w|+O(1).
\end{equation}
On the positive real axis there is a constant $C_*>0$ such that
\begin{equation}\label{eq:realgrowth}
 P(u)\sim C_* u^\kappa\exp\{A u(\log(u/(\alpha b_0))-1)\}.
\end{equation}
\end{lemma}
\begin{proof}
For the gamma factor of a character of parity $a$, write $c=1/4+a/2$. Uniform Stirling in the right half-plane gives
\[
 \log\Gamma(w/2+c)=(w/2+c-1/2)\log(w/2)-w/2+\tfrac12\log(2\pi)+O(1/|w|).
\]
The conductor term adds $(w/2+c)\log(q/\pi)$. The real part of the logarithm of $L(1/2+w,\chi)$ is uniformly bounded by the absolutely convergent logarithmic Euler product, and it tends to zero as $u\to\infty$ on the real line. The extra quadratic factor of $\xi$ contributes $2\log|w|+O(1)$ per zeta factor. Removing the central order subtracts $m\log|w|$. Summing yields \eqref{eq:stirling} with \eqref{eq:kappa}. All bounded terms have limits on the positive axis. Since $\alpha b_0=2\pi/Q_*$ and $P(u)>0$, this gives \eqref{eq:realgrowth}.
\end{proof}

\begin{lemma}[Moment upper bound]\label{lem:moment}
If \eqref{eq:moment} holds, $F$ is entire and, whenever $|u|\ge\alpha b$,
\begin{equation}\label{eq:momentupper}
 |F(u+iv)|\le M_b(\mu)\exp\{A |u|(\log(|u|/(\alpha b))-1)\}.
\end{equation}
For all $u$, $\log|F(u+iv)|\le A|u|\log(1+|u|)+O(|u|+1)$.
\end{lemma}
\begin{proof}
Uniform integrability of $e^{wt}$ and all its $w$-derivatives on compact sets follows because $e^{b e^{\alpha|t|}}$ dominates every polynomial times every ordinary exponential. Also
\[
 |F(u+iv)|\le M_b(\mu)\exp\left(\sup_{x\ge0}\{|u|x-be^{\alpha x}\}\right).
\]
The derivative vanishes at $x=\alpha^{-1}\log(|u|/(\alpha b))$ when $|u|\ge\alpha b$. Strict concavity gives the maximum and \eqref{eq:momentupper}. Below that range the supremum is $-b$.
\end{proof}

\begin{lemma}[Hadamard product and good heights]\label{lem:heights}
Assuming Dirichlet GRH, one has
\begin{equation}\label{eq:prod}
 P(w)=\prod_{\gamma>0}(1+w^2/\gamma^2)^{m_\gamma}.
\end{equation}
For every sufficiently large $R$ there is $T\in[R,R+1]$ such that uniformly for $|u|\le2$,
\begin{equation}\label{eq:good} \log|P(u\pm iT)|\ge -C T .\end{equation}
\end{lemma}
\begin{proof}
The zero count gives $\sum_{\gamma>0}m_\gamma/\gamma^2<\infty$ and the local estimate
\begin{equation}\label{eq:localcount}
 \sum_{|\gamma-y|\le1}m_\gamma\ll\log(2+|y|),
\end{equation}
where the sum is over the signed noncentral ordinates. Hadamard's theorem, order at most one and evenness give \eqref{eq:prod}: paired genus-one factors reduce to the displayed product, and the remaining exponential factor must be constant by evenness and equals one at the origin.

Remove from $[R,R+1]$ intervals of radius $c/\log(2+R)$ about all ordinates in $[R-1,R+2]$. By \eqref{eq:localcount} their total length is less than one for a sufficiently small fixed $c>0$. Choose $T$ in the remainder. Every signed ordinate then has distance at least $c/\log(2+R)$ from $T$. Taking ratios in \eqref{eq:prod} and using $|u+i(T-\gamma)|\ge|T-\gamma|$ gives
\[
 \log\frac{|P(2+iT)|}{|P(u+iT)|}
 \le\frac12\sum_{\gamma\ne0}m_\gamma\log\left(1+\frac4{(T-\gamma)^2}\right).
\]
For distances at most one the contribution is $O(\log T\log\log T)$. At distances in $[k,k+1]$ for integers $k\ge1$, use $\log(1+x)\le x$ and \eqref{eq:localcount}; the remaining sum is $O(\sum_{k\ge1}\log(2+T+k)/k^2)=O(\log T)$. Formula \eqref{eq:stirling} at $2+iT$ gives $\log|P(2+iT)|\ge-C'T$. This proves \eqref{eq:good} on the top edge; evenness gives the bottom edge. Multiplicities are included throughout; no lower bound on individual zero spacings is needed.
\end{proof}

\section{Entire division and polynomial rigidity}
\begin{lemma}[Exponential type of the quotient]\label{lem:type}
For a full-spectrum filter with $M_b(\mu)<\infty$ for any $b>0$, the quotient $E=F/P$ extends to an entire function, $E(0)=1$, of finite exponential type: $|E(w)|\le C\exp(C|w|)$.
\end{lemma}
\begin{proof}
The prescribed vanishing orders make all possible poles removable. Write $w=u+iv$. For $u\ge2$, the real part of $w\log w$ is $u\log|w|-v\arg w\ge u\log u-\frac\pi2|v|$. Lemmas \ref{lem:stirling} and \ref{lem:moment} therefore give $\log|E(w)|\le C(1+|w|)$. Use evenness of $P$ and the bound for $F$ in terms of $|u|$ when $u\le-2$.

On the portions $|u|\le2$ of the horizontal edges of the square $[-T,T]+i[-T,T]$ at a good height from Lemma \ref{lem:heights}, $F$ is bounded independently of $T$ and $|P|\ge e^{-CT}$. Thus $|E|\le e^{C'(1+T)}$ on the entire boundary. The maximum-modulus principle extends this bound to the square. For any sufficiently large $w$ take $R=|w|+1$ and a corresponding $T\le R+1$ to get the stated type estimate.
\end{proof}

\begin{lemma}[A polynomial criterion]\label{lem:poly}
Let $E$ be entire of finite exponential type. Suppose some nonnegative integer $K$ has the following property: for every $\epsilon>0$ there is a wedge angle $\theta\in(0,\pi/2)$ and a constant $C_\epsilon$ such that
\begin{equation}\label{eq:wedge}
 |E(u+iv)|\le C_\epsilon(1+|u+iv|)^K e^{\epsilon|v|}\quad\text{if }|v|\le |u|\tan\theta.
\end{equation}
Then $E$ is a polynomial.
\end{lemma}
\begin{proof}
In the upper sector $\theta\le\arg w\le\pi-\theta$ apply Phragm\'en--Lindel\"of to
\[
 E(w)e^{i\epsilon w}(w+i)^{-K}.
\]
This is holomorphic in the sector (the pole at $-i$ is outside), bounded on its rays by \eqref{eq:wedge}, and of growth at most exponential of order one. The sector opening $\pi-2\theta<\pi$ permits this order strictly below the Phragm\'en--Lindel\"of threshold. It is bounded throughout the sector. In the lower sector use $e^{-i\epsilon w}(w-i)^{-K}$. Consequently for each $\epsilon>0$,
\begin{equation}\label{eq:PW}
 |E(w)|\le C'_\epsilon(1+|w|)^K e^{\epsilon|\Ima w|}\qquad(w\in\C).
\end{equation}
By the Paley--Wiener--Schwartz theorem \cite[Thm.~7.3.1]{hormander}, the inverse Fourier transform of $E|_\R$ is a tempered distribution supported in $[-\epsilon,\epsilon]$. It is the same distribution for each $\epsilon$, hence is supported at zero. A distribution supported at one point is a finite sum of derivatives of a point mass; its Fourier transform is a polynomial. This proves the claim.

For clarity, the support conclusion can also be obtained directly by integrating $E(x)\widehat\varphi(x)$ for a smooth test function $\varphi$ supported in $[a,c]$, $a>0$, with $\widehat\varphi(w)=\int e^{-iwt}\varphi(t)\dd t$. Shift the contour to $x-iY$. Integration by parts gives $|\widehat\varphi(x-iY)|\le C_N e^{-aY}|x-iY|^{-N}$, and \eqref{eq:PW} with $\epsilon<a$ makes the shifted integral tend to zero as $Y\to\infty$. The vertical sides vanish for each fixed $Y$ by taking $N>K+2$. Negative support is handled by shifting upward. This proves support at zero without an appeal to a general support theorem.
\end{proof}

\begin{proof}[Proof of Theorem \ref{thm:main}, barrier and polynomial part]
Suppose $b\ge b_0$. First let $w=u+iv$, where $u\ge\max(2,\alpha b)$. Set $y=|v|/u$. Subtract \eqref{eq:stirling} from \eqref{eq:momentupper}, using $\alpha b_0=2\pi/Q_*$:
\begin{align}
 \log|E(w)|\le{}&-A u\log(b/b_0)+ A u\Psi(y)-\kappa\log|w|+C,\label{eq:angle}\\
 \Psi(y)={}&y\arctan y-\tfrac12\log(1+y^2).
\end{align}
Here $\Psi(0)=0$, $\Psi'(y)=\arctan y$, hence $0\le\Psi(y)\le y^2/2$. The same estimate holds with $u$ replaced by $|u|$ when $u\le-\max(2,\alpha b)$ by applying \eqref{eq:stirling} to $-w$. No symmetry of $F$ is necessary.

Choose any $\epsilon>0$ and $\delta>0$ with $A\delta/2\le\epsilon$. On $|v|\le\delta|u|$ the angular term is bounded by $A|u|\delta (|v|/|u|)/2\le\epsilon|v|$. (Equivalently, $v^2/|u|\le\delta|v|$.) The first term in \eqref{eq:angle} is nonpositive. Choose $K\ge\max(0,-\kappa)$ integral. This proves \eqref{eq:wedge}, enlarging its constant on the bounded remainder of the wedge. Lemmas \ref{lem:type} and \ref{lem:poly} now imply $E=Q$ is a polynomial.

If $b>b_0$, apply \eqref{eq:angle} with $v=0$ and let $u\to+\infty$. The exponential term forces $Q(u)\to0$, impossible for a polynomial with $Q(0)=1$. This proves \eqref{eq:main}. If $b=b_0$, the quotient is real on the real axis, and $F(u),P(u)>0$ there; hence its coefficients are real and $Q(u)>0$. Evenness follows when $\mu$ is even. The degree bound is proved next.
\end{proof}

\section{The exact moment test and attainment}
\begin{lemma}[Exact series test]\label{lem:series}
Let $\mu$ be a probability measure with entire bilateral Laplace transform $F=P Q$, where $Q$ is a real polynomial, $Q(0)=1$, and $F(u)>0$ for real $u$. For $k=\deg Q$, the admissible values of $b>0$ are exactly
\[
 (0,b_0)\ \cup\ \bigl\{b_0:\ \kappa+k<-\tfrac12\bigr\}.
\]
\end{lemma}
\begin{proof}
Since $e^{\alpha n|t|}\le e^{\alpha nt}+e^{-\alpha nt}\le2e^{\alpha n|t|}$ for $n\ge0$, monotone convergence shows that the moment is finite exactly when
\begin{equation}\label{eq:series}
 \sum_{n\ge0}\frac{b^n}{n!}\bigl(F(\alpha n)+F(-\alpha n)\bigr)<\infty.
\end{equation}
The positivity of $Q$ on $\R$ forces its degree to be even and its leading coefficient positive (including the constant case). With $P$ even, \eqref{eq:realgrowth} gives, for a constant $C'>0$,
\[
 F(\alpha n)+F(-\alpha n)\sim C' n^{\kappa+k}(n/(e b_0))^n.
\]
By Stirling for $n!$, the positive summand in \eqref{eq:series} is asymptotic to a positive constant times $(b/b_0)^n n^{\kappa+k-1/2}$. This converges for $b<b_0$, diverges for $b>b_0$, and at $b=b_0$ converges exactly when the power is less than $-1$. In particular the equality power $-1$ diverges; it cannot be accepted by a weaker real-axis growth test.
\end{proof}
This completes the degree bound in Theorem \ref{thm:main}.

\begin{proof}[Proof of Theorem \ref{thm:criterion}]
For every fixed $U$, standard polynomial bounds for $L(s,\chi)$ in bounded vertical strips (from the functional equation, absolute convergence on the right and Phragm\'en--Lindel\"of), and vertical Stirling, imply
\[
 \sup_{|u|\le U}|P(u+iv)|\ll_U (1+|v|)^{C_U} e^{-\pi d |v|/4}.
\]
Only large $|v|$ matters; the finite part is bounded by entirety. Fourier inversion defines $h$ and every derivative. Shifting the integral to $\Rea w=U$ or $-U$ is justified by the displayed bound on horizontal contour edges; with factors $(-w)^\ell$ the same proof shows that $h^{(\ell)}(t)=O_{U,\ell}(e^{-U|t|})$. Fourier inversion and the identity theorem give $\int e^{wt}h(t)\dd t=P(w)$ on all of $\C$. Conjugation and evenness make $h$ real and even, and $\int h=1$.

If an endpoint filter exists, Theorem \ref{thm:main} gives $F=PQ$; Fourier uniqueness identifies the measure with $Q(-\partial_t)h\dd t$ (integration by parts introduces $w$ per factor $-\partial_t$). Thus its continuous density is everywhere nonnegative. Conversely for a polynomial satisfying the stated degree and positivity constraints, the same integration by parts proves that its density has mass $Q(0)=1$, an entire transform $PQ$, and the prescribed zero divisor. Lemma \ref{lem:series} gives the endpoint moment. Averaging a filter with its reflection replaces $Q(w)$ by $(Q(w)+Q(-w))/2$ and cannot increase its degree; the moment is unchanged. Finally, when $h\ge0$ take $Q=1$ and use Lemma \ref{lem:series}, then invoke the universal barrier and the degree obstruction for all competitors. This proves every assertion.
\end{proof}

\begin{remark}[Additional zeros are allowed]
The argument does not restrict the extra zero divisor of $F$ at the outset. At the endpoint it proves that the extra divisor must be finite. A translated canonical kernel may keep its zero divisor but does not in general keep its critical two-sided localization coefficient. Also, when $\kappa$ is sufficiently negative, the positivity test is a finite-dimensional condition on coefficients, not a claim that every allowable polynomial is positive after application to $h$.
\end{remark}

\section{Explicit sharp examples}
\subsection{Three positive theta kernels}
For $\zeta$, the classical Riemann kernel, even in $t$, is
\[
 \Phi(t)=\sum_{n\ge1}(4\pi^2n^4e^{9t/2}-6\pi n^2e^{5t/2})e^{-\pi n^2e^{2t}},\qquad t\ge0.
\]
Its Laplace transform is $\xi(1/2+w)$ (the theta identity and two integrations by parts give this formula; see \cite[\S2.16]{titch} or \cite[\texttt{fil:lem:kernels}]{fil}). Each term is positive since $2\pi n^2e^{2t}>3$. Normalize by $\xi(1/2)>0$. The first term dominates as $t\to+\infty$, so the log of the normalized density equals $-\pi e^{2t}+O(t+1)$.

For $\chi_{-4}$ the theta kernel is
\[
 K_4(t)=2e^{3t/2}\sum_{n\ge0}(-1)^n(2n+1)e^{-\pi(2n+1)^2e^{2t}/4}.
\]
The usual primitive theta transformation (root number one) makes it even. On $t\ge0$ the summands in absolute value strictly decrease: their successive ratios are at most $3e^{-2\pi e^{2t}}<1$. The alternating sum is positive and its first term dominates at infinity, giving log density $-(\pi/4)e^{2t}+O(t+1)$. Its Laplace transform is the completed beta function at $1/2+w$, by Mellin integration on the right followed by continuation. In particular its normalizing central value is positive.

The real quadratic character modulo five has values $0,1,-1,-1,1$ on residues $0,1,2,3,4$, respectively. Its root number is one and its parity is even. Its theta kernel is
\[
 K_5(t)=2e^{t/2}\sum_{n\ge1}\chi_5(n)e^{-\pi n^2e^{2t}/5}.
\]
It is even by the primitive theta transformation (the Gauss sum is $\sqrt5$) and its Laplace transform is $\Lambda_{\chi_5}(1/2+w)$. To prove positivity for $t\ge0$, divide the tail absolute sum by the leading term and bound it by
\[
 \sum_{n\ge2}e^{-\pi(n^2-1)e^{2t}/5}
 \le\sum_{n\ge2}e^{-3\pi(n-1)/5}
 =\frac{e^{-3\pi/5}}{1-e^{-3\pi/5}}<1.
\]
The first inequality uses $n^2-1\ge3(n-1)$. In particular $K_5$ is strictly positive, has positive normalizing value at the centre, and log density $-(\pi/5)e^{2t}+O(t+1)$ as $t\to+\infty$. All three positivity arguments are unconditional. The root-number identities also follow directly from Poisson summation with the respective finite Gauss sums.

\begin{corollary}[An infinite family with exact extremum]\label{cor:examples}
Assume Dirichlet GRH. For a multiset made of $r,s,j$ copies of these three functions, the supremum of admissible $b$ among all full-spectrum filters is precisely \eqref{eq:example}. It is admissible for at least one filter exactly when $j>7r+s+2$.
\end{corollary}
\begin{proof}
There are no central zeros by the positivity just proved. Here $r_0=r$, $r_1=s$, $m=0$. Thus \eqref{eq:kappa} gives $\kappa=(7r+s-j)/4$. Convolution of the three normalized positive densities, with multiplicities $r,s,j$, is a probability density with Laplace transform $P$, by the convolution theorem. This is the canonical $h$. Apply Theorem \ref{thm:criterion}; the inequality $\kappa<-1/2$ is equivalent to the one in the assertion. When this inequality fails, no polynomial degree can repair it, even for nonsymmetric measures.
\end{proof}

\begin{center}
\begin{tabular}{@{}lccc@{}}
\toprule
Family & $\alpha$ & supremum $b_0$ & attained?\\\midrule
$\zeta$ & $2$ & $\pi$ & No\\
$L(s,\chi_{-4})$ & $2$ & $\pi/4$ & No\\
$\zeta L(s,\chi_{-4})$ & $1$ & $\pi$ & No\\
$L(s,\chi_5)$ & $2$ & $\pi/5$ & No\\
$L(s,\chi_5)^2$ & $1$ & $2\pi/5$ & No\\
$L(s,\chi_5)^3$ & $2/3$ & $3\pi/5$ & Yes\\
$\zeta L(s,\chi_5)^9$ & $1/5$ & $10\pi/5^{9/10}$ & No\\
$\zeta L(s,\chi_5)^{10}$ & $2/11$ & $11\pi/5^{10/11}$ & Yes\\\bottomrule
\end{tabular}
\end{center}

\section{Tail constants and truncation of prime averages}
For $\mu$ satisfying a moment of positive coefficient define
\[
 T_\mu(H)=\mu(\{t:|t|>H\}),\quad
 c(\mu)=\sup\{b>0:M_b(\mu)<\infty\}.
\]
Use the convention $-\log0=+\infty$. Layer-cake integration shows
\begin{equation}\label{eq:liminf}
 c(\mu)=\liminf_{H\to\infty}e^{-\alpha H}\log(1/T_\mu(H)).
\end{equation}
Indeed Markov gives the lower bound by each admissible $b$. Conversely, if the liminf is larger than $b$, for some $\eta>0$ the tail of $Y=e^{\alpha|t|}$ is eventually at most $e^{-(b+\eta)y}$, and $\mathbb E e^{bY}=1+\int_0^\infty b e^{by}\Pr(Y>y)\dd y$ is finite. This also proves the equivalence with the natural limiting tail coefficient, with no assumptions on smoothness.

\begin{theorem}[Canonical tail and sharp truncation constant]\label{thm:tail}
Let $h$ be a convolution of the normalized positive kernels in Section 5, with family size $d$. Then
\begin{equation}\label{eq:tail}
 \lim_{|t|\to\infty}-e^{-\alpha|t|}\log h(t)=b_0,\qquad
 \lim_{H\to\infty}e^{-\alpha H}\log(1/T_h(H))=b_0.
\end{equation}
For $c>0$, the smallest radius $H_{x,c}$ with $T_h(H_{x,c})\le x^{-c}$ satisfies
\begin{equation}\label{eq:quantile}
 H_{x,c}=\frac d2\log\log x+\frac d2\log(c/b_0)+o(1)\quad(x\to\infty).
\end{equation}
No other fixed full-spectrum filter with a positive moment at exponent $\alpha$ can satisfy $T_\mu(H)\le\exp(-(b_0+\eta)e^{\alpha H})$ at every sufficiently large $H$, for any $\eta>0$.
\end{theorem}
\begin{proof}
Write $h=h_1*\cdots*h_d$ and $c_i=\pi/q_i$ for the leading exponential constants of the individual kernels. They satisfy $\log h_i(t)=-c_i e^{2|t|}+O(1+|t|)$ for all $t$, by positivity and the first-term asymptotics. Thus for every $\epsilon\in(0,1/2)$,
\[
 h_i(t)\le C_\epsilon e^{-(1-\epsilon)c_i e^{2|t|}}.
\]
In the convolution integral with $\sum_i t_i=t$, AM--GM gives $\sum_i c_i e^{2|t_i|}\ge d(\prod_i c_i)^{1/d}e^{2|t|/d}=b_0 e^{\alpha|t|}$. Keep a fraction $\epsilon$ of each term to make the integrals over $t_1,\ldots,t_{d-1}$ bounded (drop the last remaining factor), obtaining $h(t)\le C'_{\epsilon}e^{-(1-2\epsilon)b_0 e^{\alpha|t|}}$.

For the lower estimate take $t>0$ large and put $t_i^*=t/d+\frac12\log(q_i/Q_*)$, so $\sum_i t_i^*=t$ and $c_i e^{2t_i^*}=b_0 e^{\alpha t}/d$. If $d\ge2$, integrate over $|t_i-t_i^*|\le\delta/d$ for $i<d$, which forces $|t_d-t_d^*|<\delta$. All coordinates are positive for sufficiently large $t$. On this box the sum of leading costs is at most $e^{2\delta}b_0 e^{\alpha t}$ and the sum of linear errors is $O(t+1)$. The box has fixed positive volume, giving $h(t)\ge C_\delta e^{-e^{2\delta}b_0 e^{\alpha t}-O(t+1)}$. For $d=1$ the individual kernel asymptotic suffices. Send $\epsilon,\delta$ to zero and use evenness to prove the first limit.

The second follows by integrating the upper bound after substituting $y=e^{\alpha t}$, and integrating the lower bound on $[H,H+\delta]$ then letting $\delta\downarrow0$. Positivity of $h$ makes $T_h$ continuous, strictly decreasing from one to zero, so inversion of the second limit gives \eqref{eq:quantile}. The last assertion follows from \eqref{eq:liminf} and Theorem \ref{thm:main}.
\end{proof}

\paragraph{Application to the explicit formula.} With $W\in C_c^\infty((1,2))$ and $S_\chi(x)=\sum_n\Lambda(n)\chi(n)W(n/x)$ (here $\Lambda(n)$ is the von Mangoldt function), the standard explicit formula under Dirichlet GRH gives a bounded normalized oscillatory error whose noncentral modes are $e^{i\gamma\log x}$. The article's full-spectrum averaging theorem (\texttt{fil:thm:globalprime}) cancels these exactly by averaging in $\log x$ against a full-spectrum filter. For the positive canonical family above, truncating the normalized average at $H$ and renormalizing changes the average of any bounded function $f$ by at most $2\|f\|_\infty T_h(H)$: split the original expectation on the retained and deleted sets and use the triangle inequality. Thus \eqref{eq:quantile} identifies the exact constant for truncating this fixed global averaging rule at error scale $x^{-c}$ in normalized bounded tests. For unnormalized explicit formulas one must also control the exponential moment weights in the main terms; the same leading double-exponential tail coefficient applies to each fixed weight $e^{vt}$ by \eqref{eq:tail}. This statement is not an optimum over arbitrary $x$-dependent filters, nor a new bound for unaveraged prime errors.

\section{Scope, further questions, and proof audit}
\paragraph{What has been resolved.} The coefficient $2\pi d/Q_*$ in the source open question is not sharp: for every fixed family a factor of two is lost by the circular averaging bound. The new constant is the exact universal scale and is the exact family-specific supremum in the examples of Section 5. Endpoint attainment there is fully classified, and general endpoint attainment reduces to a finite-dimensional positivity problem. In the arbitrary-family case the new upper bound need not, on the present proof, be the exact supremum; we have not proved that a signed canonical kernel can always be repaired without reducing its coefficient. This is a genuine remaining problem. The novelty claim is relative to the inspected repository statements, especially its explicit open question; a comprehensive priority search or peer review has not been performed.

\paragraph{Why GRH is relevant.} It identifies the complete prescribed real-frequency divisor with that of $P$; without it, cancelling the real ordinates of all zeros need not imply division by $P$. The analytic upper-bound proof would also work without GRH if one required vanishing at the full \emph{complex} zero divisor instead (with appropriate modification of the product estimate for zeros in the strip), but that is a different cancellation condition. We make no claim that these results prove Dirichlet GRH, nonabelian GRH, simplicity of zeros, or independence of zeros. No new Lean formalization was run.

\paragraph{Further research questions.}
\begin{enumerate}
 \item Is $b_0$ the supremum for every conjugation-closed family, even when its canonical signed kernel takes negative values? Determine which entire multipliers of small exponential type can enforce positivity without losing the limiting coefficient.
 \item When $\kappa<-1/2$, classify the nonnegative differential polynomials $Q(-\partial_t)h$ in the allowed degree range. Find a practical certified test of their positivity on the whole real axis. Is the canonical $h$ already positive for additional families of quadratic or complex characters?
 \item Can the method be extended to completed automorphic $L$-functions, assuming their corresponding RH and classical analytic continuation, functional equation, local zero count and Euler-product bounds? Their archimedean degree and shifts should govern an analogous gamma power and localization constant.
 \item For finite-frequency filters, can entire division by a finite divisor plus an exterior growth constraint remove the $O(\log\log B)$ excess in the article's support asymptotic? This note treats an infinite spectrum and a fixed global measure, not that finite compact-support problem.
 \item Is there a rigidity theorem for near-extremizers as $b\uparrow b_0$ under a suitable uniform moment normalization? Compactness alone is insufficient if moment constants are allowed to diverge.
\end{enumerate}

\paragraph{Audit checklist.} (i) Multiplicities, including coincident ordinates across factors, are retained in the division and in the local count. (ii) Central zeros are divided out before normalizing. (iii) The good-height construction does not require simple zeros or quantitative zero separation. (iv) The quotient's growth order is established before applying Phragm\'en--Lindel\"of. (v) The sector opening is strictly less than $\pi$; no conclusion relies on the limiting half-plane version. (vi) Both real rays are controlled, including nonsymmetric measures. (vii) The endpoint power threshold is strict, by a positive series comparison, including the harmonic boundary case. (viii) Sharpness examples use positivity proved directly, not numerical sampling. (ix) The tail extremum is explicitly scoped to fixed measures and the moment criterion. The accompanying numerical script checks representative gamma asymptotics and illustrates the series power; it is not used in the proof.

\begin{thebibliography}{9}
\bibitem{proveit} ProveIt, \emph{Generalized Riemann Hypothesis}, project README and merged article introduction, accessed 27 November 2026. \url{https://github.com/VladimirReshetnikov/ProveIt/tree/main/Analysis/GeneralizedRiemannHypothesis}
\bibitem{fil} ProveIt, \emph{Beyond the Riemann hypothesis for Dirichlet L-functions: consequences, transfers and limits}, section on positive filters, file \texttt{Article/sections/fil.tex}, blob \texttt{ee4a2a3294242cdb7594a09c5b47b350aad5c2fa}. \url{https://github.com/VladimirReshetnikov/ProveIt/blob/main/Analysis/GeneralizedRiemannHypothesis/Article/sections/fil.tex}
\bibitem{lean} \texttt{lsernau/grh-lean}, upstream commit \texttt{72fead102f663d5c7a407964a13dc0cf9d977207}, as vendored in ProveIt. \url{https://github.com/lsernau/grh-lean/tree/72fead102f663d5c7a407964a13dc0cf9d977207/grh}
\bibitem{davenport} H. Davenport, \emph{Multiplicative Number Theory}, 3rd ed., revised by H. L. Montgomery, Springer, 2000.
\bibitem{titch} E. C. Titchmarsh, \emph{The Theory of the Riemann Zeta-function}, 2nd ed., revised by D. R. Heath-Brown, Oxford, 1986.
\bibitem{zeros} M. A. Bennett, G. Martin, K. O'Bryant and A. Rechnitzer, Counting zeros of Dirichlet $L$-functions, \emph{Math. Comp.} 90 (2021), 1455--1482. \url{https://doi.org/10.1090/mcom/3599}.
\bibitem{boas} R. P. Boas, \emph{Entire Functions}, Academic Press, 1954.
\bibitem{hormander} L. H\"ormander, \emph{The Analysis of Linear Partial Differential Operators I: Distribution Theory and Fourier Analysis}, Springer, 2nd ed., 1990, Theorem 7.3.1.
\end{thebibliography}
\end{document}
